\section{Supplemental Notes}

\subsection{Note 1} \label{sup_note_1}
We experimentally determined the prior on \rb by averaging the 
	intensity of the pixels immediately outside the region of interest.
	%
	We estimated the prior on \re  by first fitting a count of $\n=1$ to the trace.
	%
	The distribution of all observed $z$-states is unimodal.
	% 
	Therefore, whichever \z{}-state is the most common, the second most common state 
	will be $\z{} \pm 1$. 
	%
	We observed that fitting $n=1$ captures this specific transition, and provides a good
	estimate of photon emission rate of a single emitter \re.

\subsection{Note 2} \label{sup_note_2}
Upon further analysis, the main limitation of \lbfcs is the estimation of the 
	intensity of a single emitter, corresponding to \re in \ours. 
	%
	\lbfcs relies on a histogram based method to determine this value, while 
	\ours is able to jointly optimize this parameter with other parameters.
	%
	Providing this value to \lbfcs led to a substantial recovery of performance.
	% 
	With \re provided, \lbfcs identified 200/300 (67\%) of traces within 1 of the true count.

\subsection{Note 3} \label{sup_note_3}
Because noise in our model is a function of intensity, quantifying the SNR of a trace 
	is not trivial. 
	%
	For simplicity, here we define SNR as the difference in intensity between 
	the first two states (\z{0} and \z{1}) divided by the standard deviation of the difference. 
	%
	In effect this is the higher bound of SNR for a given trace.

\subsection{Note 4} \label{sup_note_4}
When $\pon < \poff$, the distribution is shifted towards 0, and a majority 
	of the frames occupy the lowest two \z{}-states (\figref{fig:results:campare_kinetics}c). 
	In this regime, the true count becomes indistinguishable for all methods, without prior information. 
	%
	When $\pon \geq \poff$, this distribution is centered, or even shifted towards $n$ 
	and a larger fraction of the states are visited (\figref{fig:results:campare_kinetics}d). 
	%
	This provides ample information for the model to infer the correct $n$.

\subsection{Note 5} \label{sup_note_5}
In DNA-PAINT, the blinking rate is determined by the kinetics of single-stranded DNA binding,
which in turn depends on experimental conditions such as temperature and concentration, 
and the specific DNA sequence~\citep{jungmann_single-molecule_2010}.
	%
	As a result, the kinetic parameters \pon and \poff of our model can be 
	tuned by adjusting the temperature and imager concentration. 
	%
	As seen in (\figref{fig:results:campare_kinetics}b), \pon increases with imager concentration
	and \poff increases with temperature.
	%
	We imaged DNA origami with a known count of $\n=1$ at 25\textdegree C and an 
	imager concentration of 10 nM and measured $\pon=0.028$ and $\poff=0.072$ 
	(\figref{fig:results:campare_kinetics}b),
	%
	conditions within the poor counting 
	accuracy regime of $\pon < \poff$ (\figref{fig:results:campare_kinetics}a). % fix
	%
	To examine performance in the predicted more favorable kinetic regime,  
	we increased imager concentration and decreased temperature.
	%
	Increasing imager concentration to 30 nM raised \pon to 0.071 and 
	decreasing temperature to 13\textdegree C decreased \poff to 0.037 
	(\figref{fig:results:campare_kinetics}b).
	%
	As expected, the effects of temperature and concentration were largely independent 
	of one another \cite{jungmann_single-molecule_2010}.
	%
	We also observed a substantial decrease in SNR with decreasing temperature 
	and increasing concentration. 
	%
	This was an expected side effect of increasing concentration, because increasing 
	imager concentration increases \rb.
	%
	But the effect of temperature on SNR was surprising. 
	We hypothesize that this is due to the stabilization of partial 
	binding between imager and docker strands at low temperatures.
	%
	Balancing the increase in counting accuracy and the decrease 
	in SNR, we identified imaging conditions of 20 nM and 13\textdegree C as 
	optimal on our microscope.

\subsection{Note 6} \label{sup_note_6}
Due to the stochastic nature of blinking, multiple emitters can 
be active at any given time, which becomes increasingly likely at higher counts.
	%
	This is not compatible with the qPAINT requirement of well separated, 
	single emitter blinking events.
	%
	As a result, at higher \n the measured average dark time is longer than the true average dark time
	and results in qPAINT underestimating molecular count at higher \n
	(especially noticeable above $\n=20$, see \figref{fig:results:qpaint_counting}c). 
	%
	In \ours, the simultaneous activity of multiple emitters is modeled explicitly in both the 
	intensity and transition models (\eqref{eq:methods:p_c_given_z} and \eqref{eq:method:transition}). 
	%
	As a result, \ours avoids underestimation and can provide an accurate  
	molecular count up to $\n=30$.

\subsection{Note 7} \label{sup_note_7}
Incorporation efficiency is roughly 80 percent for each docking site \cite{strauss_2018}, 
so only a fraction of the origamis were expected to contain all 4 dockers. 
	% 
	Origamis were first imaged at 13\textdegree C with low laser power (1.5\%) and 20 nM imager concentration to 
	collect traces (4,000 frames at 200 ms exposure time) for counting with \ours (\figref{fig:results:experimental}a,b).
	%
	The system was then brought to 25\textdegree C, buffer exchange was performed, and new imager 
	was added at 10 nM.
	%
	The origamis were imaged again at high laser power (40\%),
	and post-processed with Picasso \citep{schnitzbauer_2017}, to obtain 
	super-resolution ground truth (\figref{fig:results:experimental}c).
	%
	Only origamis that had a visual count matching the designed count (1 or 4) were selected for analysis with \ours.